\section{Accessibilità}

L'accessibilità di un sito web è fondamentale per garantire che tutti gli utenti, indipendentemente dalle loro capacità fisiche, sensoriali o cognitive, possano accedere ai contenuti e ai servizi offerti. Un sito accessibile permette la fruizione delle informazioni anche a persone con disabilità, favorendo l'inclusione digitale e rispettando i principi di uguaglianza. Inoltre, l'accessibilità migliora l'usabilità generale del sito, rendendolo più semplice da navigare per tutti gli utenti, e contribuisce a soddisfare i requisiti normativi previsti dalla legge. Nel nostro caso in Italia si richiede certificazione AA.

\subsection{Separazione tra contenuto, presentazione e struttura}

In tutto il progetto è stata mantenuta una separazione tra i 3 elementi chiave di un sito web: contenuto, presentazione e struttura. In modo da migliorare l'accessibilità e la manutenibilità del sito.

\subsubsection{Assenza di CSS e JavaScript}
Il sito è completamente funzionale anche in assenza di CSS e JavaScript, garantendo così l'accesso ai contenuti anche a utenti con disabilità che utilizzano tecnologie assistive come screen reader o browser testuali.

\subsubsection{Attributi ARIA}

Sono sempre stati utilizzati gli attributi ARIA (Accessible Rich Internet Applications) per migliorare l'accessibilità delle interfacce web dinamiche. Questi attributi forniscono informazioni aggiuntive ai dispositivi di assistenza, come screen reader, aiutando gli utenti a comprendere meglio la struttura e il funzionamento del sito.

\subsubsection{Contrasti}

I colori del sito sono stati scelti per garantire un contrasto sufficiente tra testo e sfondo, in modo da facilitare la lettura anche a persone con disabilità visive. Sono stati utilizzati strumenti di verifica del contrasto per assicurarsi che i valori rispettino le linee guida WCAG. In particolare, il contrasto tra testo e sfondo deve essere almeno di 4.5:1 per il testo normale e di 3:1 per il testo grande (18pt o 14pt grassetto). Per gli elementi grafici e le icone, si è cercato di mantenere un contrasto minimo di 3:1 rispetto allo sfondo circostante. Questo sia per la modalitá scura che la modalitá chiara del sito.

\subsubsection{Lingue straniere e termini abbreviati}

Nel sito sono sempre stati utilizzati gli attributi \texttt{lang} e \texttt{abbr} quando necessari per indicare una lingua diversa dall italiano o per abbreviazioni di termini.

\subsection{Strumenti di Testing}

Per il controllo di accessibilità e presentazioni del sito sono stati utilizzati i seguenti strumenti:

\subsubsection{Total validator}

La validazione con Total Validator ha segnalato 5 falsi positivi e 1 probabile errore:
\begin{itemize}
  \item \textbf{E626 — All element names must be in lower case}. Falso positivo dovuto agli SVG inline: alcuni nomi di elementi SVG sono in camelCase (p.es.\ \texttt{linearGradient}, \texttt{feGaussianBlur}) e sono validi secondo le specifiche SVG. Il validatore applica invece la regola XHTML/HTML che impone i nomi in minuscolo, segnalando erroneamente l’\verb|<svg>|.
\item \textbf{E621 — The ATTRIBUTE NAME attribute for this tag is missing}. Nel nostro caso riguarda \texttt{xmlns} mancante su \verb|<html lang="it">|. Per documenti HTML5 serviti come \texttt{text/html} con \verb|<!DOCTYPE html>|, il namespace è implicito e \texttt{xmlns} non è richiesto (né va usato).
  \item \textbf{E627 — All attribute names must be in lower case}. Come sopra, ma per gli attributi SVG: nomi come \texttt{viewBox}, \texttt{preserveAspectRatio} o \texttt{markerUnits} sono corretti in SVG; la segnalazione nasce perché lo strumento non riconosce il namespace SVG e pretende gli attributi tutti in minuscolo.
  \item \textbf{E61 — The ATTRIBUTE NAME attribute has a boolean value, and this may be ignored or cause errors in older browsers/robots}. L’avviso riguarda gli attributi booleani (p.es.\ \texttt{required}), che in HTML5 sono validi sia in forma ridotta (\texttt{required}) sia esplicita (\texttt{required="required"}). La presenza di una delle due forme è conforme; la segnalazione è quindi un falso positivo legato alla doppia notazione prevista dallo standard.
  \item \textbf{P871 - Link text is missing}. Segnalato nel nell'\texttt{header} del sito, dove il logo è un link alla homepage: \texttt{<a href="index.php" class="site-logo" aria-label="Vai alla homepage" aria-hidden="true" tabindex="-1">}. Il testo del link non è assente, poichè è specificato nell’attributo \texttt{aria-label}.
\end{itemize}

\subsubsection{Screen reader NVDA}

É stata eseguita una lettura del sito con lo screen reader NVDA (NonVisual Desktop Access) versione 2023.4.2, in esecuzione su Windows 11. La lettura è stata effettuata con la sintesi vocale predefinita (eSpeak NG) e con la voce italiana di Microsoft.

\subsubsection{Google Lighthouse}

E stato utilizzato Google Lighthouse per verificare l'accessibilità del sito. Nelle 4 categorie fornite da Lighthouse, il sito ha ottenuto i seguenti punteggi:
\begin{itemize}
  \item Performance: 100\%
  \item Accessibility: 100\%
  \item Best Practices: 100\%
  \item SEO: 100\%
\end{itemize}

Le uniche pagine che non hanno ottenuto un punteggio di 100\% in tutte le categorie sono la pagina Contatti (a causa dell'integrazione con Google Maps) ottenendo 78\% in Best Practices e la pagina chi siamo che a causa delle immagini seppur gestite con caricamentro di webp ottiene 85\% in Performance.

\subsubsection{Ambienti di TESTING}

Il progetto è stato testato in diversi ambienti, per garantire la massima accessibilità possibile. In particolare, sono stati utilizzati i seguenti browser:

\begin{itemize}
  \item Google Chrome (versione 114.0.5735.199)
  \item Mozilla Firefox (versione 115.0.1)
  \item Microsoft Edge (versione 114.0.1823.67)
  \item Safari (versione 16.4)
\end{itemize}

Inoltre é stato eseguito sui seguenti sistemi operativi:
\begin{itemize}
  \item Windows 11
  \item macOS Ventura (versione 13.3.1)
  \item Fedora Linux (versione 42)
\end{itemize}

E sui seguenti disositivi:

\begin{itemize}
  \item Desktop
  \item Smartphone
  \item Schermi virtualizzati tramite Google Chrome
\end{itemize}